%% kapa: exam=1 round=21 year=2010 period=1 form=A subject=economics subject_ko=경제원론 seq=5 number=45
%% source: Q-Net 감정평가사 기출문제 (cst00309) artlSeq=5024181 file=extracted/1008/5024181/2010년 제21회 감정평가사 1차시험 기출문제/1교시.hwp
%% conversion: kordoc 4.24.51 -> tools/gen1.py (auto); review=proofread-sonnet; flags=none
\begin{question}{45}
X재의 생산으로 오염물질이 발생한다. X재의 수요곡선은 P = 80 - Q이고 사적 한계비용(private marginal cost, PMC)은 PMC = Q + 30이다. X재의 생산으로 사적 한계비용에 부가적으로 발생하는 사회적 한계피해액(social marginal damage, SMD)은 SMD = 2Q + 10이다. 이 경우 X재의 사회적 최적 생산량을 달성하기 위해 정부가 부과해야 하는 종량세의 크기는?
\begin{choices}
  \item 10
  \item 20
  \item 30
  \item 40
  \item 50
\end{choices}
\end{question}
\answer{3}
